% !TeX root = ../main.tex

\chapter{Introduction}\label{chapter:introduction}

\section{Motivation}

For cameras on wearables, embedded systems, and lightweight drones, useful
localisation must fit within a limited power and computation budget. Video
encoding offers an attractive resource for this purpose: H.264 has a mature,
widely deployed hardware ecosystem, and its encoders already perform motion
estimation to compress successive images~\cite{hilscher2026voca}. A geometric
estimator that can use these motion estimates could benefit from an accelerator
already available in many camera pipelines.

The appeal of dedicated hardware is concrete. One reported mobile-platform
example contrasts an \SI{80}{\milli\watt} power budget for specialized
video/audio hardware with \SI{1000}{\milli\watt} for a
CPU~\cite{chaurasia2020passthrough}.
These platform-specific budgets motivate treating video encoding as a source of
useful visual computation. Once the encoder has searched for matching blocks,
its output contains correspondence information that an odometry system could
reuse. Such correspondences can connect the two views of a stereo camera or
track scene structure across successive images. Reusing them for stereo
matching and temporal tracking is therefore a plausible route to reducing
the additional work required for localisation.

This opportunity fits naturally into a \emph{compress-then-analyze} (CTA)
architecture. A camera device can encode its observations and send the coded
streams to a companion phone for analysis~\cite{hofer2023h264cta,hilscher2026voca}.
Compression makes that communication practical, but the receiving phone still
has to infer motion from what arrives. If codec correspondences themselves
support odometry, the receiver can perform geometric estimation without first
building a pixel-based tracking frontend. The potential benefit is thus shared
across the system: dedicated coding hardware supplies motion information, the
link carries a compressed representation, and the receiver reuses that
information for localisation.

The representation also offers flexibility in where analysis is performed.
A device with spare computation could filter codec correspondences before
transmission and send the surviving observations, while a more constrained
camera could leave that work to its receiver. This makes the balance between
local computation and communication a design choice, as motivated by the
trade-off between CTA and \emph{analyze-then-compress}
(ATC)~\cite{redondi2016cta}. The choice is especially relevant when
a receiving phone must share its energy and thermal budget with other tasks.

The geometric feasibility of this idea is the prerequisite examined here:
\emph{do codec motion vectors retain enough information for useful camera
tracking without pixel-based refinement?} Their compression-oriented block
matches need not be accurate or persistent enough for this purpose. Demonstrating
a workable odometry construction would establish a reason to explore systems
built around this widely available computational resource. This thesis provides
one such construction and studies where its geometric information is sufficient.

\section{Problem Statement and Scope}

We consider calibrated stereo visual odometry using H.264 side information,
without decoded image intensities in tracking or geometric estimation, IMU
measurements, GPS, or an external pose prior. Rectification and encoding use
pixels during offline stream preparation. The frontend and estimator consume
motion-vector metadata, calibration, rectification masks, and timestamps. The
term \emph{pixel-free} applies to this estimation boundary.

The H.264 encodings use a single reference frame (\texttt{ref=1}) and disable
B-frames (\texttt{bframes=0}). In each temporal stream, inter-frame prediction
uses the preceding encoded view; stereo correspondence uses independent
I-P-P triples. This short reference structure is relevant to live systems,
where observations must be available with limited delay. It does not make
the present input preparation a live pipeline: the experiments use two-pass
encoding and reverse complete image sequences to obtain reciprocal vectors.

The restriction exposes difficulties that pixel refinement can otherwise
resolve. Codec vectors describe block-wise prediction decisions rather than
persistent scene points, their spatial support depends on the partition
structure, and their values are selected for compression rather than physical
motion~\cite{young2020fast,hilscher2026voca}. This study investigates one
construction that recovers stereo depth, validates temporal associations,
maintains identities, and integrates block-displacement observations into a
point-based estimator.

The experiments establish an analysis interface and its geometric performance.
They use prepared temporal, reverse-temporal, and shifted stereo streams;
a live hardware encoder, wireless link, and receiving phone are outside the
experimental setup. The costs of producing and transporting those inputs must
therefore be included when assessing a future distributed deployment.

\section{Research Questions}

\begin{enumerate}
  \item How much odometric accuracy remains when temporal correspondence tracking uses raw codec motion vectors without pixel-domain refinement?
  \item To what extent can block-based codec motion vectors form persistent,
        geometrically useful stereo-temporal tracks?
  \item How are the properties of the resulting tracks and input sequences related to variation in odometric accuracy?
\end{enumerate}

\section{Contributions}

Pixelless VO (PVO) provides a concrete feasibility demonstration through one
stereo pipeline. It combines codec-derived depth, persistent geometric
association, and a displacement-based observation model. Evaluation on all
eleven sequences of the established EuRoC MAV
benchmark~\cite{burri2016euroc} relates trajectory accuracy to the availability and
persistence of codec observations. The trajectory visualizations in
\autoref{sec:trajectory-geometry}, especially the easy Machine Hall sequences,
show the central result: reasonable tracking can be recovered from motion
vectors without pixel refinement. The reported PVO errors on EMH01 and EV101
provide further quantitative evidence of this geometric information.

These results support further exploration of pixel-free odometry and flexible
placement of codec-based analysis between capturing and receiving devices.
The particular stereo construction, filters, and backend are one way to test
that possibility. The measured limitations identify where alternative
constructions could improve accuracy, robustness, or the cost of preparing the
codec inputs.

\section{Thesis Structure}

Chapter~\ref{chapter:background} introduces the camera geometry, bundle-adjustment,
codec, and robust-registration concepts used throughout the thesis.
Chapter~\ref{chapter:related-work} positions the examined system with respect to
compressed-domain motion estimation, codec-aware tracking, and motion-vector
reconstruction. Chapter~\ref{chapter:method} describes the PVO compressed-domain
frontend and its integration into the visual-odometry backend.
Chapter~\ref{chapter:experimental-methodology} defines the reported configurations,
dataset, metrics, and diagnostic protocol, and Chapter~\ref{chapter:results} reports
the resulting trajectory and frontend measurements. Finally,
Chapter~\ref{chapter:discussion} interprets the operating envelope and limitations,
before Chapter~\ref{chapter:conclusion} summarizes the findings and identifies
directions for future work.
